% ============================================================================
% Informe académico — Cálculo Aplicado
% Simulación numérica de sistemas dinámicos aplicada a ciberseguridad
%
% Estructura según la Guía de Escritura de Informes Académicos:
%   Resumen, Introducción, Marco teórico, Metodología, Resultados,
%   Discusión, Conclusiones, Bibliografía.
%
% Las secciones de Resultados, la Discusión (de datos) y las Conclusiones
% quedan vacías a propósito: se completan con las simulaciones.
%
% Archivos ya generados por simulaciones.py (disponibles para descomentar).
% Cada figura existe en dos formatos: .svg (principal) y .pdf (para LaTeX;
% pdflatex no incluye SVG directamente, por eso se mantiene el PDF):
%   ../figuras/caso1_lineal.svg|.pdf        Figura caso 1 (crecimiento lineal)
%   ../figuras/caso2_exponencial.svg|.pdf   Figura caso 2 (decrecimiento exponencial)
%   ../figuras/error_vs_t.svg|.pdf          Error absoluto vs. tiempo
%   ../figuras/euler_vs_mejorado.svg|.pdf   Euler vs. Euler mejorado (h=1)
%   ../figuras/orgs_abc.svg|.pdf            Organizaciones A, B y C (h=1)
%   ../figuras/org_b_paso.svg|.pdf          Efecto del paso, organización B
%   ../figuras/discreto_hk.svg|.pdf         Modelo discreto con distintos hk
%   ../figuras/lambda_modelo.svg|.pdf       Modelo con λ (nuevas vulnerabilidades)
%   ../tablas/error.tex                     Tabla de error vs. h
%   ../tablas/mejorado.tex                  Tabla Euler vs. Euler mejorado
%   ../tablas/orgs.tex                      Tabla organizaciones (h=1)
%   ../tablas/orgs_h.tex                    Error máximo de cada organización vs. h
% ============================================================================

\documentclass[11pt,a4paper]{article}

\usepackage[T1]{fontenc}
\usepackage[utf8]{inputenc}
\usepackage[spanish]{babel}
\usepackage{amsmath,amssymb}
\usepackage{graphicx}
\usepackage{booktabs}
\usepackage{float}
\usepackage[
  colorlinks=true,
  linkcolor=black,
  urlcolor=blue,
  citecolor=blue
]{hyperref}
\usepackage[left=2.5cm,right=2.5cm,top=2.5cm,bottom=2.5cm]{geometry}

\graphicspath{{../figuras/}}

\title{\textbf{Simulación numérica de sistemas dinámicos\\aplicada a ciberseguridad}\\
\large Ecuaciones en diferencias, ecuaciones diferenciales y métodos numéricos}
\author{[Nombre y apellido]\\
\small Cálculo Aplicado — Universidad Católica del Uruguay}
\date{}

\begin{document}
\maketitle

% ---------------------------------------------------------------------------
\section*{Resumen}
% ---------------------------------------------------------------------------
% Contenido: objetivo, metodología y resultados principales.
% La conclusión principal se omitirá hasta completar la Sección de Conclusiones.

Este trabajo estudia la simulación numérica de sistemas dinámicos aplicada a
la gestión de vulnerabilidades en ciberseguridad. El objetivo es aproximar la
evolución de la cantidad de equipos vulnerables de una organización mediante
el método de Euler y su variante mejorada (Heun), analizando el efecto del
tamaño de paso $h$ sobre la precisión y el costo computacional.

Se implementaron ambos métodos en Python y se verificaron sobre dos casos con
solución de referencia conocida: crecimiento lineal ($x'=1$) y decrecimiento
exponencial ($x'=-x$, $x(t)=e^{-t}$). Sobre el segundo se midió el error
absoluto $E_n=|x(t_n)-x_n|$ para $h\in\{1,\,0{,}5,\,0{,}1,\,0{,}05,\,0{,}01\}$
y se compararon los métodos. En el modelo de remediación $V'(t)=-kV(t)$,
$V(0)=1000$, se simularon organizaciones con $k=0{,}1$, $0{,}3$ y $0{,}7$;
se estudió el modelo discreto $V_{n+1}=(1-hk)V_n$ y la extensión
$V'(t)=\lambda-kV(t)$ con nuevas vulnerabilidades.

Euler resulta exacto en el caso lineal para cualquier $h$; en el exponencial,
el error máximo cae de $3{,}68\times10^{-1}$ ($h=1$) a $1{,}85\times10^{-3}$
($h=0{,}01$), con 10 y 1000 iteraciones respectivamente. Euler mejorado
alcanza error máximo menor que Euler para igual $h$ ($6{,}6\times10^{-3}$ frente
a $1{,}9\times10^{-2}$ con $h=0{,}1$), al costo de dos evaluaciones de $f$ por
iteración. Con mayor $k$ la cantidad de equipos vulnerables decrece más
rápido; $h=1$ introduce errores de hasta $\sim200$ equipos para $k=0{,}7$,
mucho menores con $h=0{,}1$. El modelo discreto es una sucesión geométrica de
razón $1-hk$: para $hk>1$ toma valores negativos, sin sentido físico. Con
$\lambda>0$, $V(t)$ se estabiliza en $V^*=\lambda/k$.

% [Pendiente: conclusión principal, a completar junto con la Sección de
% Conclusiones.]

% ---------------------------------------------------------------------------
\section{Introducción}
% ---------------------------------------------------------------------------

Muchos problemas de la vida real involucran cantidades que cambian a lo largo
del tiempo: el crecimiento de una población, la evolución de la temperatura,
la carga de una batería o, en el ámbito de la informática, la cantidad de
equipos de una organización que permanecen expuestos a una vulnerabilidad de
seguridad. Estudiar estos fenómenos consiste, en esencia, en describir cómo
cambia una variable durante intervalos pequeños de tiempo.

Cuando no es posible (o no es necesario) resolver analíticamente una ecuación
diferencial, se puede aproximar la evolución del sistema mediante métodos
numéricos: sucesiones de valores calculadas paso a paso en una computadora.
Este proyecto parte de conceptos ya conocidos ---sucesiones, ecuaciones en
diferencias y derivadas--- para introducir progresivamente el concepto de
ecuación diferencial y su simulación computacional.

El problema elegido es la gestión de vulnerabilidades en ciberseguridad.
Supongamos que una organización posee inicialmente una determinada cantidad de
equipos vulnerables y que, a medida que se aplican actualizaciones y parches,
dicha cantidad disminuye. Nos interesa estudiar cómo evoluciona ese proceso y
cómo influyen la rapidez de remediación, el tamaño del paso utilizado en la
simulación, el método numérico empleado, las condiciones iniciales y la
aparición de nuevas vulnerabilidades.

\textbf{Objetivo general.}
Estudiar métodos numéricos sencillos para aproximar la evolución de sistemas
dinámicos y analizar su aplicación en un problema relacionado con la gestión de
vulnerabilidades informáticas.

\textbf{Objetivos específicos.}
\begin{itemize}
  \item Relacionar ecuaciones en diferencias con ecuaciones diferenciales.
  \item Implementar el método de Euler y una mejora del mismo.
  \item Analizar el efecto del tamaño del paso y estudiar el error numérico.
  \item Comparar métodos en términos de precisión y costo computacional.
  \item Aplicar los métodos a un modelo simplificado de ciberseguridad e
        interpretar los resultados.
  \item Investigar conceptos vinculados con gestión de vulnerabilidades, CVE,
        CVSS y MITRE ATT\&CK.
\end{itemize}

El informe está organizado como sigue: en la Sección~\ref{sec:marco} se
presenta el marco teórico (de las ecuaciones en diferencias a los métodos
numéricos y a los conceptos de ciberseguridad); en la
Sección~\ref{sec:metodologia} se describe la metodología de implementación y
los experimentos planificados; la Sección~\ref{sec:resultados} recoge los
resultados numéricos; la Sección~\ref{sec:discusion} discute los hallazgos, y
la Sección~\ref{sec:conclusiones} resume las conclusiones.

% ---------------------------------------------------------------------------
\section{Marco teórico}
\label{sec:marco}
% ---------------------------------------------------------------------------

\subsection{De las ecuaciones en diferencias a las ecuaciones diferenciales}
\label{sec:marco-diferencias}

Consideremos una variable $x(t)$ que cambia con el tiempo. La derivada
\begin{equation}
  x'(t)
\end{equation}
representa la tasa de cambio instantánea de $x$. Para un incremento pequeño
$h$, la definición de derivada permite utilizar la aproximación
\begin{equation}
  x'(t) \approx \frac{x(t+h) - x(t)}{h}.
  \label{eq:aprox-derivada}
\end{equation}

Si una ecuación diferencial tiene la forma
\begin{equation}
  x'(t) = f(t, x),
  \label{eq:edo}
\end{equation}
entonces, sustituyendo en~\eqref{eq:aprox-derivada},
\begin{equation}
  \frac{x(t+h) - x(t)}{h} \approx f(t, x).
\end{equation}
Despejando $x(t+h)$,
\begin{equation}
  x(t+h) \approx x(t) + h\, f(t, x).
  \label{eq:euler-continuo}
\end{equation}

Si definimos $t_n = t_0 + nh$ e identificamos $x_n \approx x(t_n)$,
la expresión~\eqref{eq:euler-continuo} se convierte en la \emph{ecuación en
diferencias}
\begin{equation}
  x_{n+1} = x_n + h\, f(t_n, x_n),
  \label{eq:euler-discreto}
\end{equation}
que constituye la idea básica del método de Euler. Obsérvese que una ecuación
diferencial se transforma así en una sucesión discreta: el valor en el paso
siguiente se obtiene a partir del valor actual, sin necesidad de resolver la
ecuación analíticamente.

\subsection{Método de Euler}
\label{sec:marco-euler}

El método de Euler aproxima la solución del problema de valor inicial
\begin{equation}
  x'(t) = f(t, x), \qquad x(t_0) = x_0.
\end{equation}
Construye la sucesión
\begin{equation}
  \boxed{\; x_{n+1} = x_n + h\, f(t_n, x_n) \;},
  \qquad t_{n+1} = t_n + h,
  \label{eq:euler}
\end{equation}
donde el parámetro $h > 0$ se denomina \emph{tamaño del paso}. Geométricamente,
en cada paso se avanza desde el punto $(t_n, x_n)$ en la dirección de la
pendiente $f(t_n, x_n)$, evaluada al comienzo del intervalo.

\subsection{Método de Euler mejorado}
\label{sec:marco-mejorado}

El método de Euler utiliza únicamente la pendiente al comienzo de cada
intervalo. Una posible mejora (método de Heun o Euler mejorado) consiste en
realizar primero una \emph{predicción},
\begin{equation}
  \tilde{x}_{n+1} = x_n + h\, f(t_n, x_n),
\end{equation}
y utilizar luego el \emph{promedio} entre la pendiente inicial y la pendiente
predicha al final del intervalo:
\begin{equation}
  \boxed{\;
    x_{n+1} = x_n + \frac{h}{2}\Bigl[f(t_n, x_n) + f\bigl(t_{n+1},
    \tilde{x}_{n+1}\bigr)\Bigr].
  \;}
  \label{eq:euler-mejorado}
\end{equation}
Cada iteración requiere dos evaluaciones de $f$, en lugar de una. En la
literatura numérica, Euler es un método de orden~1 y Euler mejorado es un
método de orden~2: al reducir $h$ a la mitad, el error máximo del primero se
reduce aproximadamente a la mitad y el del segundo aproximadamente a una
cuarta parte~\cite{burden2016}.

\subsection{Error numérico}
\label{sec:marco-error}

Para comparar una aproximación numérica $x_n$ con una solución de referencia
$x(t_n)$ se define el \emph{error absoluto}
\begin{equation}
  E_n = \bigl| x(t_n) - x_n \bigr|.
  \label{eq:error}
\end{equation}
Algunos problemas admiten solución de referencia conocida (por ejemplo,
$x(t) = e^{-t}$ para $x' = -x$, $x(0)=1$), lo que permite medir cuantitativamente
la calidad de la aproximación sin requerir resolución analítica de la ecuación
diferencial en el informe. Del error se derivan medidas útiles como el
\emph{error final} $E_N$ (en el último tiempo simulado) y el \emph{error
máximo} $E_{\max} = \max_n E_n$ sobre toda la trayectoria.

\subsection{Gestión de vulnerabilidades}
\label{sec:marco-gestion}

Una \emph{vulnerabilidad} es una debilidad en un sistema informático que puede
ser explotada por un atacante para comprometer su confidencialidad,
integridad o disponibilidad~\cite{nistvuln}. La \emph{gestión de
vulnerabilidades} es el proceso cíclico de identificar, evaluar, tratar y
reportar estas debilidades de forma continua~\cite{nist80040}. El tratamiento
habitual consiste en aplicar un \emph{parche de seguridad}: una actualización
de software que corrige la vulnerabilidad. El \emph{tiempo de corrección o
remediación} es el lapso transcurrido entre la aparición (o detección) de la
vulnerabilidad y la aplicación efectiva del parche en el activo afectado.

Para estandarizar la identificación y la priorización existen dos referentes
internacionales:

\begin{itemize}
  \item \textbf{CVE} (\emph{Common Vulnerabilities and Exposures}): catálogo
        público, mantenido por MITRE, que asigna un identificador estándar
        (con formato \texttt{CVE-AAAA-NNNNN}) a cada vulnerabilidad conocida
        y pública, evitando la ambigüedad entre proveedores~\cite{cve}.
  \item \textbf{CVSS} (\emph{Common Vulnerability Scoring System}): sistema
        de puntuación mantenido por FIRST que cuantifica la severidad de una
        vulnerabilidad en una escala de 0 a 10, a partir de métricas de
        explotabilidad e impacto~\cite{cvss}. La puntuación CVSS es una de
        las herramientas de \emph{priorización} más utilizadas: a igualdad de
        otras condiciones, primero se parchean los sistemas con vulnerabilidades
        de mayor severidad.
\end{itemize}

Mantener sistemas sin actualizar tiene consecuencias directas: un equipo
vulnerable permanece expuesto durante todo el tiempo de remediación, lo que
amplía la \emph{superficie de ataque} de la organización---el conjunto de
puntos de entrada potencialmente explotables. Desde el punto de vista
operativo, la rapidez con la que se apagan vulnerabilidades se puede modelar
como un parámetro de remediación $k$: un valor de $k$ mayor indica una
respuesta más rápida (más equipos parcheados por unidad de tiempo).

\subsection{MITRE ATT\&CK: técnica T1190}
\label{sec:marco-t1190}

MITRE ATT\&CK for Enterprise es una base de conocimiento que documenta las
tácticas y técnicas que los adversarios utilizan en ataques reales~\cite{mitre}.
La técnica \textbf{T1190 --- Exploit Public-Facing Application} consiste en
explotar una vulnerabilidad de una aplicación expuesta a Internet (servidor
web, servicio de correo, VPN, panel de administración, etc.) para obtener un
acceso inicial a la red de la víctima~\cite{t1190}.

Esta técnica se relaciona directamente con la existencia de sistemas que
permanecen vulnerables: una aplicación pública con una vulnerabilidad sin
parchear es, precisamente, el objetivo de T1190. Reducir el tiempo de
remediación disminuye la ventana durante la cual ese objetivo existe y, por
tanto, reduce la superficie de ataque. En términos del modelo estudiado en
este informe, el parámetro $k$ representa la capacidad de respuesta del
proceso de remediación: a mayor $k$, más rápido decrece la cantidad de
equipos vulnerables $V(t)$ y menos tiempo queda expuesta la organización a
técnicas como T1190.

% ---------------------------------------------------------------------------
\section{Metodología}
\label{sec:metodologia}
% ---------------------------------------------------------------------------

\subsection{Implementación}

Los métodos numéricos se implementaron en Python en los módulos
\texttt{euler.py} (funciones \texttt{euler} y \texttt{euler\_mejorado}) y
\texttt{simulaciones.py} (experimentos y generación de figuras y tablas).
Ambas funciones reciben la función $f(t,x)$, la condición inicial $x_0$, los
tiempos inicial $t_0$ y final $tf$, y el tamaño de paso $h$; devuelven los
valores de tiempo $t_n$, los valores aproximados $x_n$ y la cantidad de pasos
realizados.

El método de Euler se obtuvo de la ecuación en
diferencias~\eqref{eq:euler}. El Euler mejorado implementa el esquema de
predicción-corrección de la ecuación~\eqref{eq:euler-mejorado}: en cada paso
se calcula primero $\tilde{x}_{n+1}$ (pendiente inicial) y luego se promedia
con $f(t_{n+1}, \tilde{x}_{n+1})$ (pendiente predicha).

En todos los experimentos, la simulación comienza en $t_0 = 0$ y avanza con
paso constante $h$ mientras $t_n \le tf$ (se incluye el instante final cuando
$h$ divide a $tf$; en caso contrario, el último punto es el inmediatamente
anterior).

\subsection{Experimentos}

Se planificaron los siguientes experimentos, en línea con los objetivos
específicos:

\begin{enumerate}
  \item \textbf{Caso 1 --- crecimiento lineal.}
        $x'(t) = 1$, $x(0) = 0$, con solución de referencia $x(t) = t$,
        para $0 \le t \le 10$ y $h \in \{1,\, 0{,}5,\, 0{,}1\}$.
        Verifica la implementación en un caso donde la aproximación debe ser
        exacta (la pendiente es constante).
  \item \textbf{Caso 2 --- decrecimiento exponencial.}
        $x'(t) = -x$, $x(0) = 1$, con solución de referencia $x(t) = e^{-t}$,
        para $0 \le t \le 10$ y $h \in \{1,\, 0{,}5,\, 0{,}1,\, 0{,}01\}$.
  \item \textbf{Error y tamaño del paso.}
        Para el caso 2, se calcula el error absoluto
        $E_n = |x(t_n) - x_n|$ con $h \in \{1,\, 0{,}5,\, 0{,}1,\, 0{,}05,\,
        0{,}01\}$; se grafica $E_n$ contra $t$ y se registran número de pasos,
        error final y error máximo.
  \item \textbf{Comparación de métodos.}
        Se aplican Euler y Euler mejorado al caso 2 con
        $h \in \{1,\, 0{,}5,\, 0{,}1\}$ y se comparan valor final, error final,
        error máximo y cantidad de pasos.
  \item \textbf{Aplicación: gestión de vulnerabilidades.}
        Modelo $V'(t) = -k V(t)$, $V(0) = V_0 = 1000$, con solución de
        referencia $V(t) = V_0 e^{-kt}$, para tres organizaciones
        (A: $k = 0{,}1$; B: $k = 0{,}3$; C: $k = 0{,}7$),
        $0 \le t \le 20$ y $h \in \{1,\, 0{,}5,\, 0{,}1\}$.
        Se registran los valores de $V$ en $t = 5, 10, 15$ y el error frente
        a la referencia.
  \item \textbf{Modelo discreto.}
        Análisis de la ecuación en diferencias
        $V_{n+1} = (1 - hk)\, V_n$ obtenida al aplicar Euler al modelo
        anterior: tipo de sucesión, expresión explícita, límite y valores de
        $hk$ con sentido físico.
  \item \textbf{Modelo con nuevas vulnerabilidades.}
        $V'(t) = \lambda - k V(t)$ con $V_0 = 1000$, $k = 0{,}3$ y
        $\lambda \in \{0,\, 50,\, 100,\, 200\}$, simulado con Euler.
\end{enumerate}

\subsection{Criterios de análisis}

\begin{itemize}
  \item \textbf{Exactitud:} se mide con el error absoluto
        $E_n$~\eqref{eq:error} frente a la solución de referencia, reportando
        error final y error máximo.
  \item \textbf{Efecto del paso:} se compara el error y el número de pasos
        al variar $h$; el costo computacional se asocia a la cantidad de
        iteraciones (y, en Euler mejorado, a las dos evaluaciones de $f$ por
        iteración).
  \item \textbf{Sentido físico del modelo discreto:} se examina cuándo
        $(1-hk)$ produce sucesiones no negativas, condición necesaria para
        que $V_n$ represente una cantidad de equipos.
  \item \textbf{Interpretación organizacional:} los resultados se contrastan
        con los conceptos de gestión de vulnerabilidades, CVE, CVSS y técnicas
        ATT\&CK de la Sección~\ref{sec:marco-gestion}.
\end{itemize}

% ---------------------------------------------------------------------------
\section{Resultados}
\label{sec:resultados}
% ---------------------------------------------------------------------------
% SECCIÓN A COMPLETAR. Presentar aquí los datos obtenidos de forma objetiva,
% con figuras y tablas numeradas y referenciadas desde el texto, sin
% interpretación (la interpretación va en Discusión).
%
% Los archivos de figuras y tablas ya fueron generados por simulaciones.py
% (ver listado al comienzo de este archivo). Para incorporarlos, basta con
% descomentar los bloques correspondientes y ajustar números de figura/tabla,
% textos de pie y referencias en el cuerpo.

\subsection{Verificación: casos 1 y 2}

% \begin{figure}[H]
%   \centering
%   \includegraphics[width=0.85\textwidth]{caso1_lineal.pdf}
%   \caption{Caso 1: crecimiento lineal. Solución de referencia $x(t)=t$
%   comparada con las aproximaciones de Euler para $h=1$, $h=0{,}5$ y
%   $h=0{,}1$, en $0 \le t \le 10$.}
%   \label{fig:caso1}
% \end{figure}

% \begin{figure}[H]
%   \centering
%   \includegraphics[width=0.85\textwidth]{caso2_exponencial.pdf}
%   \caption{Caso 2: decrecimiento exponencial. Solución de referencia
%   $x(t)=e^{-t}$ comparada con las aproximaciones de Euler para
%   $h=1$, $h=0{,}5$, $h=0{,}1$ y $h=0{,}01$, en $0 \le t \le 10$.}
%   \label{fig:caso2}
% \end{figure}

% Describir objetivamente lo observado (p.~ej., en el caso 1 la
% aproximación coincide con la referencia para cualquier $h$; en el caso 2
% las curvas se acercan a la referencia a medida que $h$ disminuye).

\subsection{Error numérico y tamaño del paso}

% \begin{figure}[H]
%   \centering
%   \includegraphics[width=0.85\textwidth]{error_vs_t.pdf}
%   \caption{Error absoluto $E_n = |x(t_n) - x_n|$ en función del tiempo para
%   el caso 2 ($x'=-x$, $x(0)=1$), con $h = 1$, $0{,}5$, $0{,}1$, $0{,}05$ y
%   $0{,}01$.}
%   \label{fig:error}
% \end{figure}

% \input{../tablas/error}

% Completar con las observaciones directas sobre el error (sin interpretar).

\subsection{Comparación entre Euler y Euler mejorado}

% \begin{figure}[H]
%   \centering
%   \includegraphics[width=0.85\textwidth]{euler_vs_mejorado.pdf}
%   \caption{Euler y Euler mejorado con $h=1$ frente a la solución de
%   referencia $x(t)=e^{-t}$, en $0 \le t \le 10$.}
%   \label{fig:mejorado}
% \end{figure}

% \input{../tablas/mejorado}

\subsection{Gestión de vulnerabilidades: organizaciones A, B y C}

% \begin{figure}[H]
%   \centering
%   \includegraphics[width=0.85\textwidth]{orgs_abc.pdf}
%   \caption{Evolución de equipos vulnerables $V(t)$ para las organizaciones
%   A ($k=0{,}1$), B ($k=0{,}3$) y C ($k=0{,}7$), con $V_0=1000$. Curvas
%   continuas: solución de referencia $V(t)=1000\,e^{-kt}$; puntos: Euler con
%   $h=1$.}
%   \label{fig:orgs}
% \end{figure}

% \begin{figure}[H]
%   \centering
%   \includegraphics[width=0.85\textwidth]{org_b_paso.pdf}
%   \caption{Organización B ($k=0{,}3$): solución de referencia y
%   aproximaciones de Euler con $h=1$, $h=0{,}5$ y $h=0{,}1$, en
%   $0 \le t \le 20$.}
%   \label{fig:org-b}
% \end{figure}

% \input{../tablas/orgs}

% \input{../tablas/orgs_h}

\subsection{Modelo discreto}

% \begin{figure}[H]
%   \centering
%   \includegraphics[width=0.85\textwidth]{discreto_hk.pdf}
%   \caption{Modelo discreto $V_{n+1}=(1-hk)V_n$ con $k=0{,}3$ y $V_0=1000$
%   para $hk=0{,}3$ (positivo), $hk=1{,}5$ (alterna el signo) y
%   $hk=2{,}5$ (diverge).}
%   \label{fig:discreto}
% \end{figure}

% Registrar la sucesión obtenida, valores de $V_n$ para distintos $hk$.

\subsection{Modelo con aparición de nuevas vulnerabilidades}

% \begin{figure}[H]
%   \centering
%   \includegraphics[width=0.85\textwidth]{lambda_modelo.pdf}
%   \caption{Modelo $V'(t) = \lambda - 0{,}3\,V(t)$ con $V_0=1000$, para
%   $\lambda = 0$, $50$, $100$ y $200$. Marcas: simulación con Euler
%   ($h=0{,}1$); curvas: solución de referencia. Niveles estacionarios
%   $V^* = \lambda/k$.}
%   \label{fig:lambda}
% \end{figure}

% Describir el comportamiento observado en cada caso de $\lambda$.

% ---------------------------------------------------------------------------
\section{Discusión}
\label{sec:discusion}
% ---------------------------------------------------------------------------
% SECCIÓN A COMPLETAR (interpretación de los resultados de la Sección 4).
% Sugerencias de contenido:
%   - qué sucede con el error y con el número de iteraciones al disminuir h;
%     compromiso entre precisión y costo computacional;
%   - cuál de los dos métodos presenta menor error y si depende de h;
%     operaciones por iteración; cuándo conviene un método más preciso;
%   - qué representa k, cuál organización remedia más rápido, cómo cambia la
%     gráfica al aumentar k; factores reales que justifican valores distintos
%     de k entre organizaciones;
%   - análisis de hk > 1 y valores negativos en el modelo discreto;
%   - comportamiento del modelo con λ (decrecer, crecer o estabilizarse) y
%     su interpretación organizacional;
%   - limitaciones del modelo (equipos idénticos, k constante, no considera
%     ataques reales ni la aparición de vulnerabilidades ---salvo en el modelo
%     con λ---, etc.) y qué aspectos incorporaría un modelo más realista;
%   - análisis global: precisión ↔ tamaño del paso ↔ costo computacional,
%     estabilidad y significado informático de los resultados.

% ---------------------------------------------------------------------------
\section{Conclusiones}
\label{sec:conclusiones}
% ---------------------------------------------------------------------------
% SECCIÓN VACÍA A PROPÓSITO. Debe contener:
%   - síntesis de los hallazgos más importantes;
%   - respuesta a los objetivos planteados en la Introducción;
%   - implicancias o recomendaciones;
%   - posibles líneas futuras.
% No presentar resultados nuevos.

% ---------------------------------------------------------------------------
\section{Bibliografía}
% ---------------------------------------------------------------------------
% Formato APA. Las referencias siguientes corresponden a las fuentes técnicas
% citadas en el marco teórico; completar con las que se utilicen en la
% Discusión.

\begin{thebibliography}{99}
\bibitem{burden2016}
Burden, R.~L., \& Faires, J.~D. (2016).
\emph{Numerical Analysis} (10th~ed.). Cengage Learning.

\bibitem{cve}
MITRE. (s.\,f.).
\emph{CVE --- Common Vulnerabilities and Exposures}.
https://www.cve.org/

\bibitem{cvss}
Forum of Incident Response and Security Teams (FIRST). (s.\,f.).
\emph{Common Vulnerability Scoring System}.
https://www.first.org/cvss/

\bibitem{nistvuln}
National Institute of Standards and Technology (NIST). (s.\,f.).
\emph{Vulnerabilities}.
Glossary. https://csrc.nist.gov/glossary/vulnerability

\bibitem{nist80040}
National Institute of Standards and Technology (NIST). (2022).
\emph{Guide to Enterprise Patch Management Planning: Preventive Maintenance
for Technology} (NIST Special Publication 800-40 Revision 4).
https://doi.org/10.6028/NIST.SP.800-40r4

\bibitem{mitre}
MITRE. (s.\,f.).
\emph{MITRE ATT\&CK for Enterprise}.
https://attack.mitre.org/

\bibitem{t1190}
MITRE. (s.\,f.).
\emph{T1190: Exploit Public-Facing Application}.
ATT\&CK for Enterprise.
https://attack.mitre.org/techniques/T1190/
\end{thebibliography}

\end{document}
